\chapter{置换群与不可约表示的 Young 图构造}
\label{chap:permutation-group}

有限群里最特殊的一类是对称群 $S_n$（$n$ 个对象的全部置换）。它特殊有两个理由。其一，Cayley 定理告诉我们任何一个 $n$ 阶有限群都同构地嵌进某个 $S_m$，所以 $S_n$ 的表示论是一切有限群表示论的「母本」：把 $S_n$ 弄清楚，任意有限群都能看成某个 $S_m$ 的子群，再通过限制表示拉回来。其二，它直接服务于物理：$n$ 个全同粒子交换标签，交换动作本身就是 $S_n$ 的元素——这一条下一章\chapref{chap:identical-particles}会专门展开。

本章的问题只有一个：\textbf{$S_n$ 有哪些不可约表示？怎样把它们一个个构造出来、算出维数和矩阵元？} 我们沿三步走：先把置换本身讲透（循环分解、共轭类、类大小），再用 Young 图与 Young 对称化子把每个分拆对应到一个不可约表示，最后用钩长公式与 Young--Yamanouchi 规则把维数和相邻对换的矩阵元真正算出来。$S_3$、$S_4$ 会反复手算，作为每一步的检验。

\section{置换、循环结构与共轭类}
\label{sec:sn-cycle-conjugacy}

置换群研究的是“标签怎样互换”。如果有 $n$ 个彼此可区分的位置 $1,2,\dots,n$，一个置换 $\pi$ 就是一一映射
\[
\pi:\{1,2,\dots,n\}\to\{1,2,\dots,n\}.
\]
例如循环 $(123)$ 表示 $1\mapsto2$、$2\mapsto3$、$3\mapsto1$，而没有写出的标签保持不动。本文将乘积 $\pi\sigma$ 定义为先施行 $\sigma$、再施行 $\pi$；因此 $(12)(23)$ 必须逐个追踪标签，不能把两个括号交换。

为什么物理需要这种看似离散的标签游戏？对全同粒子来说，粒子编号本身没有物理意义。把电子 1 与电子 2 的标签交换以后，哈密顿量不应改变。于是 $S_n$ 自然作用在 $n$ 粒子 Hilbert 空间上，而玻色/费米统计就是对这种置换表示的特殊要求。

设
\[
X_n=\{1,2,\dots,n\}.
\]
所有双射 $\pi:X_n\to X_n$ 在函数复合下构成群 $S_n$。本讲义继续采用“右边先作用”的乘法约定，
\[
(\sigma\pi)(i)=\sigma(\pi(i)).
\]
阶数为
\[
|S_n|=n!,
\]
因为一个置换等价于把 $n$ 个不同元素排成一个有序序列。

\begin{definition}[循环与循环类型]
若置换 $\pi$ 依次把
\[
a_1\mapsto a_2,\quad a_2\mapsto a_3,\quad\dots,\quad a_k\mapsto a_1,
\]
并固定其余元素，则记为
\[
(a_1a_2\cdots a_k),
\]
称为一个 $k$-循环。若两个循环涉及的数字集合不相交，则称它们不交。

若一个置换的不交循环分解中恰有 $m_k$ 个 $k$-循环，则称
\[
1^{m_1}2^{m_2}\cdots n^{m_n}
\]
为该置换的循环类型，其中
\begin{equation}
\sum_{k=1}^n k m_k=n.
\label{eq:partition-cycle-type}
\end{equation}
\label{def:cycle-type}
\end{definition}

任意置换都能唯一地分解成不交循环的乘积，唯一性只差各循环的先后顺序和每个循环内部的循环移位。这个结论可以直接从轨道理解：由子群 $\langle\pi\rangle$ 作用在 $X_n$ 上，每个轨道恰好形成一个循环，不同轨道互不相交。

\begin{proposition}[不交循环分解的唯一性]
任意 $\pi\in S_n$ 都可写成若干两两不交循环的乘积；若忽略 $1$-循环，则这种分解在循环因子次序以及循环内部循环移位之外唯一。

\end{proposition}

\begin{proof}
固定 $a\in X_n$，考虑序列
\[
a,\pi(a),\pi^2(a),\dots.
\]
因为 $X_n$ 有限，必有最小正整数 $r$ 使 $\pi^r(a)=a$。于是轨道
\[
\{a,\pi(a),\dots,\pi^{r-1}(a)\}
\]
给出循环
\[
(a\ \pi(a)\ \cdots\ \pi^{r-1}(a)).
\]
对不在该轨道中的元素重复此过程，有限步后覆盖 $X_n$，不同轨道天然不交。

若有另一组不交循环分解，则包含 $a$ 的循环必须精确包含 $a$ 在 $\langle\pi\rangle$ 下的轨道，否则不能实现同一个映射。因此每个循环对应的支持集合被 $\pi$ 唯一决定，循环内部的次序也由 $a\mapsto\pi(a)$ 唯一决定，只允许改变起点。故唯一性成立。
\end{proof}

不交循环彼此交换，因此置换的许多性质只由循环长度决定。例如一个 $k$-循环的阶是 $k$，从而若
\[
\pi=c_1c_2\cdots c_r
\]
是不交循环分解，循环长度分别为 $\ell_1,\dots,\ell_r$，则
\begin{equation}
\operatorname{ord}(\pi)=\operatorname{lcm}(\ell_1,\dots,\ell_r).
\label{eq:permutation-order-lcm}
\end{equation}

\begin{workedexample}
证明两个不交循环交换。

\textbf{解.} 设 $c,d$ 支持集合不交。对任意 $x$：若 $x$ 不在两个支持中，两者都固定；若只在 $c$ 支持中，$d$ 固定 $x$ 及 $c(x)$，故 $cd(x)=c(x)=dc(x)$；若只在 $d$ 支持中同理。因此 $cd$ 与 $dc$ 对每个 $x$ 作用相同。
\end{workedexample}


二循环 $(ab)$ 称为对换。任意 $k$-循环可写成 $k-1$ 个对换的乘积，例如
\begin{equation}
(a_1a_2\cdots a_k)
=(a_1a_k)(a_1a_{k-1})\cdots(a_1a_2).
\label{eq:cycle-as-transpositions}
\end{equation}
因此任意置换都能写成对换的乘积。不同分解所需的对换数并不唯一，但奇偶性却是唯一的。

\begin{theorem}[置换奇偶性的良定义性]
若一个置换 $\pi\in S_n$ 可写成 $r$ 个对换的乘积，也可写成 $s$ 个对换的乘积，则
\[
r\equiv s\pmod 2.
\]
因此可以定义符号同态
\begin{equation}
\operatorname{sgn}:S_n\to\{\pm1\},
\qquad
\operatorname{sgn}(\pi)=(-1)^r,
\label{eq:sign-homomorphism}
\end{equation}
其中 $r$ 是任意一个对换分解的长度。

\end{theorem}

\begin{proof}
考虑多项式
\[
\Delta(x_1,\dots,x_n)=\prod_{1\le i<j\le n}(x_i-x_j),
\]
即 Vandermonde 反对称多项式。置换 $\pi$ 作用在变量标签上，定义
\[
(\pi\Delta)(x_1,\dots,x_n)
=\Delta(x_{\pi^{-1}(1)},\dots,x_{\pi^{-1}(n)}).
\]
一个对换只交换两个变量，因此使 $\Delta$ 变号。若 $\pi$ 是 $r$ 个对换的乘积，就有
\[
\pi\Delta=(-1)^r\Delta.
\]
若另一个分解长度为 $s$，同一个线性作用又必须满足
\[
\pi\Delta=(-1)^s\Delta.
\]
由于 $\Delta\neq0$，只能有 $(-1)^r=(-1)^s$，故 $r,s$ 同奇偶。

乘积 $\sigma\pi$ 的对换分解可以直接串接，因此
\[
\operatorname{sgn}(\sigma\pi)=\operatorname{sgn}(\sigma)\operatorname{sgn}(\pi),
\]
故 \eqnrefs{eq:sign-homomorphism} 是群同态。
\end{proof}

符号同态的核
\[
A_n=\operatorname{Ker}(\operatorname{sgn})
\]
称为交错群，由全部偶置换组成，是 $S_n$ 的指数二正规子群。这个例子把\chapref{chap:group-basics}的同态核理论和本章联系了起来。

有了循环分解和奇偶性，现在推导 $S_n$ 的共轭类分类：两个置换共轭当且仅当它们有相同的循环类型，因此共轭类与 $n$ 的整数分拆一一对应。

对一般有限群，找共轭类往往是主要困难；对 $S_n$，答案极其整齐。共轭的几何意义只是“重新命名被置换的标签”，所以它不会改变循环长度。

\begin{theorem}[$S_n$ 的共轭类分类]
两个置换 $\sigma,\pi\in S_n$ 共轭，当且仅当它们具有相同的循环类型。因而 $S_n$ 的共轭类与 $n$ 的整数分拆一一对应。

\end{theorem}

\begin{proof}
先设
\[
\pi=(a_1\cdots a_k)
\]
是一个循环。任取 $\tau\in S_n$，直接作用于任意标签即可验证
\begin{equation}
\tau(a_1\cdots a_k)\tau^{-1}
=(\tau(a_1)\cdots\tau(a_k)).
\label{eq:cycle-conjugation}
\end{equation}
因此共轭只重新命名循环中的数字，不改变循环长度。对不交循环乘积逐个应用 \eqnrefs{eq:cycle-conjugation}，可知共轭置换循环类型相同。

反过来，若
\[
\pi=(a_1\cdots a_{k_1})(b_1\cdots b_{k_2})\cdots,
\]
\[
\sigma=(a'_1\cdots a'_{k_1})(b'_1\cdots b'_{k_2})\cdots
\]
具有相同循环类型，就定义一个置换 $\tau$ 把每个 $a_i\mapsto a'_i$、$b_i\mapsto b'_i$，并对所有循环逐一对应。它是 $X_n$ 的双射。由 \eqnrefs{eq:cycle-conjugation}，$\tau\pi\tau^{-1}=\sigma$。故充分性也成立。
\end{proof}

这一定理的核心公式是
\[
\tau(a_1a_2\cdots a_m)\tau^{-1}
=(\tau(a_1)\tau(a_2)\cdots\tau(a_m)).
\]
也就是说，共轭只是在循环记号里把每个标签统一重新命名，并不会改变循环长度。因此循环类型在共轭下不变。

反过来，若两个置换拥有完全相同的循环长度列表，就可以把第一个置换每个循环里的标签依次送到第二个对应循环里的标签，从而构造一个双射 $\tau$。上式随后保证第二个置换等于 $\tau\sigma\tau^{-1}$。所以“同一共轭类”与“同一循环类型”完全等价。

所谓整数分拆，就是把 $n$ 写成非增正整数之和
\[
\lambda=(\lambda_1,\lambda_2,\dots,\lambda_r),
\qquad
\lambda_1\ge\lambda_2\ge\cdots\ge\lambda_r>0,
\qquad
\sum_i\lambda_i=n.
\]
它与循环类型 $1^{m_1}2^{m_2}\cdots$ 是同一信息的两种写法：$\lambda_i$ 直接列出每个循环的长度，$m_k$ 则记录长度为 $k$ 的循环有多少个。

\begin{proposition}[给定循环类型的共轭类大小]
循环类型为
\[
1^{m_1}2^{m_2}\cdots n^{m_n},
\qquad \sum_k km_k=n
\]
的共轭类大小为
\begin{equation}
|C_{(m_1,\dots,m_n)}|
=\frac{n!}{\displaystyle\prod_{k=1}^n k^{m_k}m_k!}.
\label{eq:sn-class-size}
\end{equation}

\end{proposition}

\begin{proof}
先把 $n$ 个数字按某种顺序填入所有循环位置，总共有 $n!$ 种填法。但同一个 $k$-循环
\[
(a_1a_2\cdots a_k)
=
(a_2\cdots a_ka_1)
\]
有 $k$ 种循环移位写法，所以每个 $k$-循环带来一个 $k$ 倍重复；共有 $m_k$ 个这样的循环，给出 $k^{m_k}$。另外，同长度的 $m_k$ 个互不相交循环彼此交换，交换它们的先后顺序不产生新置换，又有 $m_k!$ 倍重复。对所有 $k$ 相乘，得到分母
\[
\prod_k k^{m_k}m_k!,
\]
于是得到 \eqnrefs{eq:sn-class-size}。
\end{proof}

\begin{example}[{$S_4$ 的五个共轭类}]
$4$ 的分拆为
\[
4,
\quad 3+1,
\quad 2+2,
\quad 2+1+1,
\quad1+1+1+1.
\]
对应循环类型依次是 $4^1$、$3^11^1$、$2^2$、$2^11^2$、$1^4$。用 \eqnrefs{eq:sn-class-size} 得类大小
\[
6,\quad 8,\quad3,\quad6,\quad1,
\]
总和为 $24=|S_4|$。因此 $S_4$ 恰有五个不等价复不可约表示。下一节将看到，它们同样由五个 Young 图标记。
\end{example}

\begin{workedexample}
求 $S_5$ 中循环类型 $2^11^3$ 与 $3^11^2$ 的类大小。

\textbf{解.} 用 \eqnrefs{eq:sn-class-size}：
\[
|C_{2,1^3}|=\frac{5!}{2^1 1!\,1^3 3!}=10,
\]
\[
|C_{3,1^2}|=\frac{5!}{3^1 1!\,1^2 2!}=20.
\]
\end{workedexample}


\section{Young 图、Young 盘与不可约表示}
\label{sec:young-representation}

前面已经知道 $S_n$ 的共轭类由循环类型决定，而循环类型恰好由整数 $n$ 的分拆标记。例如
\[
4=4=3+1=2+2=2+1+1=1+1+1+1
\]
对应 $S_4$ 的五个共轭类。第二章又告诉我们，有限群的不等价复不可约表示个数等于共轭类个数。因此“整数分拆”不仅分类群元，也应该能够给不可约表示贴标签。Young 图就是把一个分拆画成方格图，从而把这种一一对应变成可操作的组合工具。

但 Young 图本身只告诉我们“有哪一种对称类型”，还没有给出表示空间。为了构造具体空间，需要把粒子标签 $1,\dots,n$ 填进方格得到 Young 盘，然后规定同一行做对称化、同一列做反对称化。行对称与列反对称的结合，正好制造出一种不可再分的混合置换对称性。Young 对称算符和 Specht 模就是把这句话精确化。

\begin{definition}[Young 图与共轭分拆]
设 $\lambda=(\lambda_1,\dots,\lambda_r)\vdash n$。Young 图 $[\lambda]$ 由 $n$ 个方格组成，第 $i$ 行有 $\lambda_i$ 个方格，各行左对齐。把图沿主对角线转置得到的 Young 图对应分拆 $\lambda'$，称为共轭分拆。
\label{def:young-diagram}
\end{definition}

例如 $n=4$ 的五个 Young 图分别对应
\[
[4],\quad[3,1],\quad[2,2],\quad[2,1,1],\quad[1,1,1,1].
\]
这五个图既与五个共轭类一一对应，也将与五个不可约表示一一对应，但这两个“一一对应”不是说某个循环类直接对应某个同形 Young 图；它们只是恰好由同一个分拆集合分别标记。

$S_4$ 的 Young 图如图 \figref{fig:s4-young-diagrams} 所示。
\begin{figure}[htbp]
\centering
\begin{tikzfig}[scale=.95]
\foreach \x/\rows/\lab in {0/{4}/{$[4]$},3/{3,1}/{$[3,1]$},6/{2,2}/{$[2,2]$},9/{2,1,1}/{$[2,1,1]$},12/{1,1,1,1}/{$[1^4]$}}{
  \node[lecturelabel] at (\x+1.0,-2.3) {\lab};
}
% [4]
\foreach \i in {0,...,3}{\draw (\i*.48,0) rectangle ++(.48,-.48);}
% [3,1]
\foreach \i in {0,...,2}{\draw (3+\i*.48,0) rectangle ++(.48,-.48);} \draw (3,-.48) rectangle ++(.48,-.48);
% [2,2]
\foreach \r in {0,1}{\foreach \i in {0,1}{\draw (6+\i*.48,-\r*.48) rectangle ++(.48,-.48);}}
% [2,1,1]
\foreach \i in {0,1}{\draw (9+\i*.48,0) rectangle ++(.48,-.48);} \draw (9,-.48) rectangle ++(.48,-.48); \draw (9,-.96) rectangle ++(.48,-.48);
% [1^4]
\foreach \r in {0,...,3}{\draw (12,-\r*.48) rectangle ++(.48,-.48);}
\end{tikzfig}
\caption{$S_4$ 的五个 Young 图，对应 $4$ 的五个整数分拆。}
\label{fig:s4-young-diagrams}
\end{figure}

\begin{definition}[Young 盘与标准 Young 盘]
把数字 $1,2,\dots,n$ 各填入 Young 图的一个方格，得到 Young 盘 $T$。若每一行从左到右严格递增、每一列从上到下严格递增，则称为标准 Young 盘（标准 Young tableau）。
\label{def:young-tableau}
\end{definition}

同一个形状有 $n!$ 种带标签填法，但绝大多数彼此包含大量冗余。标准盘的数目最终正好等于对应不可约表示的维数。

整数分拆给出 Young 图，但 Young 图本身只是几何对象。要把它变成表示论的工具，需要在张量空间上构造 Young 算符：行方向对称化、列方向反对称化。

先不要急着看公式。若两个标签位于 Young 盘同一行，我们希望交换它们时状态保持不变；若两个标签位于同一列，我们希望交换它们时状态变号。于是对一整行里的所有置换求和会产生“完全对称化”，而对一整列的所有置换按奇偶性带符号求和会产生“完全反对称化”。

形式化：设 $R_T$ 是只在各行内部重新排列数字的置换集合，$C_T$ 是只在各列内部重新排列数字的置换集合。它们都是 $S_n$ 的子群，且
\[
R_T\cap C_T=\{e\},
\]
因为一个非平凡置换不可能同时只在行内和只在列内移动数字。下面的求和不是普通数字相加，而是在\textbf{群代数 $\CC[S_n]$}中把每个置换当作一组基向量做线性组合（$\CC[S_n]$ 的元素就是 $\sum_{\pi}a_\pi\,\pi$，加法、乘法按置换复合展开）。在这个群代数中定义三个元素：
\begin{align}
a_T&=\sum_{r\in R_T}r,\label{eq:row-symmetrizer}\\
b_T&=\sum_{c\in C_T}\operatorname{sgn}(c)c,\label{eq:column-antisymmetrizer}\\
e_T&=a_Tb_T.\label{eq:young-symmetrizer}
\end{align}
$a_T$ 负责行对称，$b_T$ 负责列反对称；二者相乘 $e_T=a_Tb_T$ 就是\textbf{Young 对称算符}。它虽然通常还不是归一化投影，但平方后只差一个常数，因此本质上就是投影到一个不可约对称性块。关键不是“画格子”，而是让图的行与列分别编码两种相反的置换对称性。

对任意 $r_0\in R_T$，由重排定理
\[
r_0a_T=\sum_{r\in R_T}r_0r=a_T.
\]
对任意 $c_0\in C_T$，
\[
b_Tc_0=\operatorname{sgn}(c_0)b_T.
\]
于是 $e_T=a_Tb_T$ 同时把“行对称、列反对称”编码进一个群代数元素。

\begin{example}[{{$S_3$ 的形状 $\left[2,1\right]$}}]
取标准盘
\[
T=\begin{array}{|c|c|}\hline 1&2\\\hline 3\\\cline{1-1}\end{array}.
\]
则
\[
R_T=\{e,(12)\},
\qquad
C_T=\{e,(13)\}.
\]
因此
\[
a_T=e+(12),
\qquad
b_T=e-(13),
\]
\[
e_T=(e+(12))(e-(13))
=e+(12)-(13)-(132).
\]
这个显式展开同时展示了每一项怎样来自“先行对称，再列反对称”的结构。
\end{example}

\begin{workedexample}
对 $S_3$ 的盘
\[
T=\begin{array}{|c|c|}\hline1&2\\\hline3\\\cline{1-1}\end{array}
\]
计算 $a_T,b_T,e_T$。

\textbf{解.}
\[
R_T=\{e,(12)\},\qquad C_T=\{e,(13)\},
\]
故
\[
a_T=e+(12),\qquad b_T=e-(13),
\]
\[
e_T=e+(12)-(13)-(132).
\]
\end{workedexample}


如果只看群代数中的 $(12)$、$(13)$ 仍然觉得抽象，可以让它真正作用在三粒子张量积上。写一个未对称化乘积态
\[
\Psi(1,2,3)=\phi_a(1)\phi_b(2)\phi_c(3).
\]
把置换看作粒子标签的重排，Young 算符给出
\begin{align*}
e_T\Psi
&=[e+(12)-(13)-(132)]\Psi\\
&=\phi_a(1)\phi_b(2)\phi_c(3)
 +\phi_a(2)\phi_b(1)\phi_c(3)\\
&\quad-\phi_a(3)\phi_b(2)\phi_c(1)
 -\phi_a(2)\phi_b(3)\phi_c(1).
\end{align*}
现在可以直接检查：左作用交换盘中同一行的标签 $1,2$ 不改变这一组合。需要区分算符次序：$e_T=a_Tb_T$ 先由 $b_T$ 对列反对称化，再由 $a_T$ 对行对称化；后一步一般不保持所有左列交换的反对称性。精确规律是 $r e_T=e_T$ 与 $e_T c=\operatorname{sgn}(c)e_T$，分别是左行作用和右列作用，不能把后者误读成 $c e_T=\operatorname{sgn}(c)e_T$。

\begin{workedexample}
证明反对称化算符作用在两个相同单粒子态上给零。

\textbf{解.} 对 $n=2$，
\[
\mathcal A=\frac12(I-P_{12}).
\]
若态为 $\phi\otimes\phi$，则 $P_{12}(\phi\otimes\phi)=\phi\otimes\phi$，所以
\[
\mathcal A(\phi\otimes\phi)=0.
\]
一般 $n$ 中只要两个列向量相同，Slater 行列式为零。
\end{workedexample}


投影算符满足 $P^2=P$。Young 算符通常没有恰好归一化成投影，但它满足“平方等于自身的常数倍”，这就是本质幂等性。

\begin{theorem}[Young 算符的本质幂等性]
对 Young 盘 $T$，存在只依赖其形状 $\lambda$ 的非零常数 $c_\lambda$，使
\begin{equation}
e_T^2=c_\lambda e_T.
\label{eq:young-quasi-idempotent}
\end{equation}
因此
\[
p_T=\frac{1}{c_\lambda}e_T
\]
是幂等元。常数 $c_\lambda$ 的具体数值暂不预设；在得到 Specht 模维数与钩长公式以后，我们会反过来算出它恰好等于所有钩长的乘积。

\end{theorem}

\begin{proof}
先注意 $e_T\neq0$：在展开
\[
e_T=\sum_{r\in R_T}\sum_{c\in C_T}\operatorname{sgn}(c)rc
\]
时，不同的 $(r,c)$ 不会给出同一个乘积。若 $rc=r'c'$，则
\[
r^{-1}r'=cc'^{-1}\in R_T\cap C_T=\{e\},
\]
所以 $r=r',c=c'$。特别地，单位元只由 $r=c=e$ 产生，其系数为 $1$，不会在展开中抵消。

接着考察任意 $x\in\mathbb C[S_n]$ 对应的 $e_Txe_T$。对 $r_0\in R_T$，重排求和指标给出 $r_0a_T=a_T$；对 $c_0\in C_T$，有 $b_Tc_0=\operatorname{sgn}(c_0)b_T$。所以
\[
r_0(e_Txe_T)c_0
=\operatorname{sgn}(c_0)e_Txe_T.
\]
这些变换律本身还没有排除其他双陪集上的系数。关键是 Young 行列冲突：若 $\pi\notin R_TC_T$，则存在对换 $\tau\in R_T$ 和对换 $c\in C_T$，使 $\tau\pi=\pi c$。对满足上述变换律的 $z=\sum_\pi a_\pi\pi$，比较这一基向量的系数便有
\[
a_{\tau\pi}=a_\pi,\qquad a_{\pi c}=-a_\pi,
\]
从而 $a_\pi=0$。因此支撑只能留在 $R_TC_T$；在这一个双陪集中，分解 $rc$ 唯一，且系数由单位元系数决定：
\[
a_{rc}=\operatorname{sgn}(c)a_e,
\qquad z=a_ea_Tb_T=a_e e_T.
\]
行列冲突与分解唯一性的完整组合证明见\chapref{app:young-proofs}的 Young 一维性引理（\eqnrefs{eq:young-bi-equivariance}）。
因此存在标量函数 $\lambda_T(x)$ 使
\[
e_Txe_T=\lambda_T(x)e_T.
\]
取 $x=e$ 得 $e_T^2=c_Te_T$。若 $T'=\pi T$ 与 $T$ 同形，则 $e_{T'}=\pi e_T\pi^{-1}$，故
\[
e_{T'}^2=\pi e_T^2\pi^{-1}=c_Te_{T'},
\]
所以常数只依赖形状，可记为 $c_\lambda$。还需排除 $c_\lambda=0$。在群代数 $A=\mathbb C[S_n]$ 上考虑右乘算符
\[
R_{e_T}:A\to A,\qquad x\mapsto xe_T.
\]
在群元基 $\{g\}_{g\in S_n}$ 中，$R_{e_T}g=ge_T$ 的 $g$ 分量等于 $e_T$ 的单位元系数，而该系数刚才已经证明为 $1$。因此
\[
\Tr R_{e_T}=n!\neq0.
\]
若 $c_\lambda=0$，则 $e_T^2=0$，从而 $R_{e_T}^2=0$；幂零算符的迹必须为零，与上式矛盾。故 $c_\lambda\neq0$，于是 $p_T=c_\lambda^{-1}e_T$ 确为幂等元。
\end{proof}

由此定义左理想
\begin{equation}
S^\lambda=\mathbb C[S_n]e_T.
\label{eq:specht-module}
\end{equation}
它称为形状 $\lambda$ 的 Specht 模。不同盘但相同形状只是在群代数中被共轭置换重新标记，因此给出等价表示；不同形状则给出不等价表示。

\begin{theorem}[Young--Specht 分类定理]
在复数域上，$S_n$ 的不等价不可约表示与 $n$ 的整数分拆 $\lambda\vdash n$ 一一对应。对任意形状为 $\lambda$ 的 Young 盘 $T$，Specht 模
\[
S^\lambda=\mathbb C[S_n]e_T
\]
是一个不可约表示空间；同形状不同 Young 盘得到等价表示，不同形状得到不等价表示。其维数将在下面的标准盘基定理中确定。

\end{theorem}

\begin{proof}
下面先在正文中给出分类定理的完整逻辑链；其中涉及行列冲突与本原幂等元的一组技术引理，在\chapref{app:young-proofs} 中逐式证明。

首先，若 $T'=\pi T$，则行群和列群分别满足
\[
R_{T'}=\pi R_T\pi^{-1},
\qquad
C_{T'}=\pi C_T\pi^{-1},
\]
从而
\[
e_{T'}=\pi e_T\pi^{-1}.
\]
由 $e_{T'}=\pi e_T\pi^{-1}$ 可知
\[
\mathbb C[S_n]e_{T'}=\mathbb C[S_n]e_T\pi^{-1}.
\]
因此右乘 $\pi^{-1}$ 定义
\[
\Phi_\pi(xe_T)=xe_T\pi^{-1}=x\pi^{-1}e_{T'},
\]
它与左正则作用交换，并由右乘 $\pi$ 给出逆映射，所以同形状表示等价。

其次，若两个形状 $\lambda,\mu$ 不同，就按行长的字典序把较大的记为 $\lambda$。Young 行列冲突引理说明：对形状 $\lambda$ 的任意盘与形状 $\mu$ 的任意盘，必能找到两个数字在前者同行、在后者同列；于是相应的“列反对称子 $\times$ 行对称子”必为零。对任意置换重标号仍成立，因此一个方向上的全部交织映射都被迫为零。若两个 Specht 模等价，则其同构的逆映射会给出该方向上的非零交织映射，矛盾，所以不同形状给出不等价表示。

\noindent\emph{手算看一眼冲突在哪里。} 取 $\lambda=[3,1]$、$\mu=[2,2]$，各写一个标准盘：
\[
T_\lambda=\begin{array}{|c|c|c|}\hline 1&2&3\\ \hline 4\\ \cline{1-1}\end{array},
\qquad
T_\mu=\begin{array}{|c|c|}\hline 1&2\\ \hline 3&4\\ \hline\end{array}.
\]
找一对数在 $T_\lambda$ 同行、在 $T_\mu$ 同列。先看 $1$ 与 $3$：在 $T_\lambda$ 里同在第一行，在 $T_\mu$ 里 $1$ 与 $3$ 都在第一列——\textbf{正是这一对冲突}。它在 $a_{T_\lambda}$ 下被对称化、在 $b_{T_\mu}$ 下被反对称化，$a_{T_\lambda}b_{T_\mu}$ 作用时这一对既对称又反对称，只能是 $0$。（核对其余对：$2,4$ 在 $T_\lambda$ 不同行，$1,2$ 在 $T_\mu$ 不同列，都不作数。）正是这种「某对格子在前者同行、后者同列」的必然性，让不同形状之间的交织映射全灭。\chapref{app:young-proofs} 把这一论证对任意 $\lambda>\mu$ 逐步展开。

第三，幂等元 $p_T=c_\lambda^{-1}e_T$ 对应左正则表示中的一个不变子空间。若这个子空间还能进一步约化，则 $p_T$ 可以分裂成两个非零正交幂等元；但 Young 基本引理给出
\[
p_T x p_T=\lambda(x)p_T
\]
对所有 $x\in\mathbb C[S_n]$ 成立，这正是本原幂等元判据，因此不能再分，故 $S^\lambda$ 不可约。

最后，分拆数 $p(n)$ 等于 $S_n$ 共轭类数，也等于复不可约表示等价类数。上面已经为每个分拆构造出一个两两不等价不可约表示，因此这些表示已经穷尽全部不可约表示。
\end{proof}

\begin{workedexample}
列出 $S_5$ 的全部共轭类类型，并验证类数等于 $5$ 的分拆数 $p(5)=7$。

\textbf{解.} 分拆为
\[
5,
4+1,
3+2,
3+1+1,
2+2+1,
2+1+1+1,
1^5.
\]
对应七种循环类型，故七个共轭类，也有七个不等价复不可约表示。
\end{workedexample}


\begin{theorem}[标准 Young 盘基定理]
设 $\lambda\vdash n$，记 $f^\lambda$ 为形状 $\lambda$ 的标准 Young 盘个数。则 Specht 模满足
\begin{equation}
\dim S^\lambda=f^\lambda.
\label{eq:specht-dimension-standard-tableaux}
\end{equation}
更具体地，固定一个形状为 $\lambda$ 的参考盘 $T_0$ 后，可以选取由标准盘标记的一组向量作为 $S^\lambda$ 的基。

\end{theorem}

\begin{proof}
先在经典 tabloid 模型中证明标准基定理。把只记每一行包含哪些数字、而忽略同行内部次序的对象记为 $\{T\}$；这些 tabloids 张成置换模 $M^\lambda$。对每个盘 $T$ 定义列反对称的 polytabloid
\[
\varepsilon_T
=\sum_{c\in C_T}\operatorname{sgn}(c)\{cT\}.
\]
若盘先在每一行内按递增顺序排列，而某一列仍出现逆序，就在第一次逆序处构造 Garnir belt。这是第一步\textbf{张成}：Garnir 关系把 $\varepsilon_T$ 表成若干个严格更高序的 polytabloids；由于同形盘只有有限多个，反复使用该关系必在有限步后终止，最终只留下行、列都递增的标准盘。因此标准 polytabloids 张成经典 Specht 模。举个最小的例子：$[2,1]$ 形只能填出两个标准盘
\[
T_1=\begin{array}{|c|c|}\hline 1&2\\ \hline 3\\ \cline{1-1}\end{array},
\qquad
T_2=\begin{array}{|c|c|}\hline 1&3\\ \hline 2\\ \cline{1-1}\end{array},
\]
任何非标准填法（例如 $1,3$ 在第一行、$2$ 在第二行但列上不成递增）都会被 Garnir 关系化到这两个中的某一个——这正是 $f^{[2,1]}=2$ 的先声。

线性无关性需要一个比“首项不同”更精确的三角性。对 tabloid $X$ 定义
\[
N_{k,r}(X)
=\#\{1,\dots,k\}\text{ 中位于前 }r\text{ 行的数字数}.
\]
若 $T$ 标准，则每一列中的小数字已经尽可能靠上，所以任意列置换 $c\in C_T$ 都满足
\[
N_{k,r}(\{cT\})\le N_{k,r}(\{T\})
\qquad \forall k,r,
\]
且 $c\neq e$ 时至少有一处严格不等。按照这些整数定义的支配偏序排列标准 tabloids 后，$\varepsilon_T$ 在 $\{T\}$ 上的系数为 $1$，而其中可能出现的其他标准 tabloids 全部严格位于它的一侧。于是标准 polytabloids 到标准 tabloids 的坐标矩阵是单位三角矩阵，因此它们线性无关。经典 Specht 模的维数正好是 $f^\lambda$。

最后把这件事与正文采用的群代数左理想 $S^\lambda=Aa_Tb_T$ 联系起来。令
\[
\omega(g)=\operatorname{sgn}(g)g,
\qquad g\in S_n,
\]
并线性延拓为群代数自同构。若 $T^{\mathsf T}$ 是转置盘，则
\[
\omega(a_T)=b_{T^{\mathsf T}},
\qquad
\omega(b_T)=a_{T^{\mathsf T}},
\]
因此
\[
\omega(a_Tb_T)=b_{T^{\mathsf T}}a_{T^{\mathsf T}},
\]
右端正是经典 tabloid 构造对应的 Young 对称化子。映射 $\omega(g)=\operatorname{sgn}(g)g$ 只是给每个群元乘一个符号、不改变向量空间维数，它的作用是把「行对称化」与「列反对称化」互换（转置盘 $T^{\mathsf T}$ 正好把行变成列）。用 $\omega$ 扭曲模作用不改变向量空间维数，而转置又给出标准 $\lambda$-盘与标准 $\lambda^{\mathsf T}$-盘的一一对应，所以
\[
\dim Aa_Tb_T=f^{\lambda^{\mathsf T}}=f^\lambda.
\]
Garnir 关系和上述三角性在\chapref{app:young-proofs} 中逐式展开，从而得到 \eqnrefs{eq:specht-dimension-standard-tableaux}。
\end{proof}

Specht 模给出了 $S_n$ 的全部不可约表示，但它们的维数需要另行计算。钩长公式提供一个纯组合的闭式表达式。

上一理论告诉我们“标准 Young 盘个数 = 不可约表示维数”。逐个枚举标准盘在 $n$ 大时很慢。钩长公式把维数变成一个局部乘积。

\begin{theorem}[钩长公式]
设 $\lambda\vdash n$，对每个方格 $u$ 定义钩长
\[
h(u)=1+(u\text{ 右侧方格数})+(u\text{ 下方方格数}).
\]
则
\begin{equation}
f^\lambda
=\frac{n!}{\displaystyle\prod_{u\in\lambda}h(u)}.
\label{eq:hook-length-formula}
\end{equation}

\end{theorem}

\begin{proof}
标准 Young 盘中的最大数字 $n$ 必须放在图的一个可移去角格，否则它右边或下边还存在格子就会违反递增性。删去 $n$ 所在角格，得到某个 $\mu\vdash n-1$ 的标准盘；反之，对任意可移去角格，把 $n$ 放回去都会得到标准盘。因此
\begin{equation}
f^\lambda=\sum_{\mu\nearrow\lambda}f^\mu,
\label{eq:young-branching-count}
\end{equation}
其中求和遍历所有从 $\lambda$ 删除一个角格得到的 $\mu$。

对 $n$ 作归纳。假设所有 $n-1$ 阶图满足钩长公式，把
\[
f^\mu=\frac{(n-1)!}{\prod_{v\in\mu}h_\mu(v)}
\]
代入 \eqnrefs{eq:young-branching-count}。删除一个角格只会使它左侧同一行与上方同一列的钩长各减少 $1$。把各项通分后，得到组合恒等式
\[
\sum_{\mu\nearrow\lambda}
\frac{1}{\prod_{v\in\mu}h_\mu(v)}
=\frac{n}{\prod_{u\in\lambda}h_\lambda(u)},
\]
从而推出 \eqnrefs{eq:hook-length-formula}。\chapref{app:young-proofs} 给出这个通分恒等式的一种逐因子证明。

\noindent\emph{拿 $[2,1]$ 手算这个通分恒等式。} 主图 $\lambda=[2,1]$ 的钩长为：左上格钩住「自己＋右格＋下格」共 $3$，右上格钩长 $1$，下格钩长 $1$，故 $\prod_{u\in\lambda}h_\lambda(u)=3\cdot1\cdot1=3$，右端 $n/\prod h=3/3=1$。$\lambda$ 有两个可移去角格：删右上格得到 $\mu_1=[2]$，钩长 $2,1$，$\prod=2$，$1/\prod=1/2$；删下格得到 $\mu_2=[1,1]$，钩长 $2,1$，$\prod=2$，$1/\prod=1/2$。左端 $\sum_{\mu\nearrow\lambda}1/\prod h_\mu=\tfrac12+\tfrac12=1$，正好等于右端。这就直观看到「删一个角格让左行、上列钩长各减 $1$」如何把两个子项凑成整图的 $n/\prod h$。代入归纳即得 $f^{[2,1]}=3!/(3\cdot1\cdot1)=2$。
\end{proof}

钩长公式看起来像一个神奇乘积，最好先理解每个钩是什么。对 Young 图中的一个方格 $u$，它的钩包含该格本身、同一行右边所有格以及同一列下方所有格；钩长度 $h(u)$ 就是这些格子的总数。例如形状 $(3,2)$ 左上角方格的钩长度为 $4$：它自己、右边两个格、下面一个格。

公式
\[
f^\lambda=\frac{n!}{\prod_{u\in\lambda}h(u)}
\]
给出标准 Young 盘数，也就是对应 Specht 模的维数。分子 $n!$ 是把 $1,\dots,n$ 任意填入 $n$ 个格子的总排列数；分母的钩乘积精确扣除了“必须沿行和列递增”带来的约束。正文证明通过分支递推建立这个结果，而不是把分母解释成简单的独立除法因子；上面的计数图像只用于帮助记忆结构。

\begin{workedexample}
证明共轭 Young 图对应的 Specht 模维数相同。

\textbf{解.} 共轭图把行与列互换，但钩长只是“右边格数 + 下边格数 +1”，转置后右、下互换，所有钩长的多重集合不变。因此由钩长公式
\[
f^{\lambda'}=f^\lambda.
\]
\end{workedexample}


\begin{corollary}[Young 算符的标准归一化]
前面本质幂等关系中的常数为
\begin{equation}
c_\lambda=\prod_{u\in\lambda}h(u).
\label{eq:hook-product}
\end{equation}
因此标准归一化的本原幂等元可以写成
\[
p_T=\frac{e_T}{\displaystyle\prod_{u\in\lambda}h(u)}.
\]

\end{corollary}

\begin{proof}
考虑群代数 $A=\mathbb C[S_n]$ 上的右乘算符
\[
R_{p_T}:A\to A,
\qquad x\longmapsto xp_T.
\]
因为 $p_T^2=p_T$，它是投影，像空间正是 $Ap_T=S^\lambda$，故
\[
\Tr R_{p_T}=\dim S^\lambda=f^\lambda.
\]
在群元基 $\{g\}_{g\in S_n}$ 中，$R_{p_T}g=gp_T$ 的 $g$ 分量总等于 $p_T$ 的单位元系数。$e_T$ 的单位元系数为 $1$，所以 $p_T=c_\lambda^{-1}e_T$ 的单位元系数为 $c_\lambda^{-1}$，从而
\[
\Tr R_{p_T}=\frac{n!}{c_\lambda}.
\]
结合标准盘基定理与钩长公式
\[
f^\lambda=\frac{n!}{\prod_{u\in\lambda}h(u)}
\]
立即得到 \eqnrefs{eq:hook-product}。这也说明钩乘积并不是先验塞进 Young 算符的归一化常数，而是由“群代数投影的迹 = 不可约空间维数”强制确定的。
\end{proof}

\begin{example}[{$S_4$ 的不可约表示维数}]
对 $[4]$，钩长为 $4,3,2,1$，故
\[
f^{[4]}=\frac{4!}{4\cdot3\cdot2\cdot1}=1.
\]
对 $[3,1]$，钩长为 $4,2,1,1$，故
\[
f^{[3,1]}=\frac{24}{8}=3.
\]
对 $[2,2]$，钩长为 $3,2,2,1$，故
\[
f^{[2,2]}=\frac{24}{12}=2.
\]
共轭图 $[2,1,1]$ 与 $[1^4]$ 的维数分别为 $3,1$。于是
\[
1^2+3^2+2^2+3^2+1^2=24=|S_4|,
\]
与正则表示维数平方和公式完全吻合。
\end{example}

\begin{workedexample}
对 Young 图 $[3,2]$ 计算每个方格的钩长，并求不可约表示维数。

\textbf{解.} Hook 长度为
\[
\begin{array}{|c|c|c|}\hline
4&3&1\\\hline
2&1\\\cline{1-2}
\end{array}
\]
所以乘积为 $24$，
\[
f^{[3,2]}=\frac{5!}{24}=5.
\]
\end{workedexample}


\begin{workedexample}
求 $S_5$ 全部不可约表示维数，并验证平方和为 $120$。

\textbf{解.} 对七个分拆用钩长公式得到
\[
1,4,5,6,5,4,1.
\]
于是
\[
1+16+25+36+25+16+1=120=5!.
\]
\end{workedexample}


Young 方法到这里已经给出一般结构，但第一次接触时最好先手算一个最小例子。对 $S_3$，分拆只有
\[
(3),\qquad(2,1),\qquad(1,1,1).
\]
前两个极端形状分别意味着完全行对称和完全列反对称，因此必然给出平凡表示与符号表示。中间形状 $(2,1)$ 既要求部分对称又要求部分反对称，所以不可能是一维；钩长公式给出
\[
f^{(2,1)}=\frac{3!}{3\cdot1\cdot1}=2.
\]
维数平方和 $1^2+2^2+1^2=6=|S_3|$，说明三种表示已经完备。这样 Young 图、钩长度和 Burnside 维数平方和第一次在同一个例子里闭合。

Specht 模给出了不可约表示的存在性，但还没有给出显式矩阵。Young--Yamanouchi 基提供了一组\emph{正交基}，使得每个相邻对换 $(i,i+1)$ 的表示矩阵有极简单的有理函数形式。这在实际计算 $S_n$ 表示矩阵时极为方便，也是量子化学中组态相互作用程序的基础工具。

\begin{definition}[轴向距离]
设 $T$ 是标准 Young 盘，数字 $i$ 位于第 $r_i$ 行第 $c_i$ 列。记其 content 为 $q_T(i)=c_i-r_i$，并定义从 $i$ 到 $j$ 的\textbf{轴向距离}为
\begin{equation}
d_T(i,j)=q_T(j)-q_T(i)=c_j-r_j-c_i+r_i.
\label{eq:axial-distance}
\end{equation}
若相邻数字 $i,i+1$ 在同一行，则 $d_T(i,i+1)=1$；若在同一列，则 $d_T(i,i+1)=-1$。
\label{def:axial-distance}
\end{definition}

\begin{theorem}[Young--Yamanouchi 正交表示]
对形状 $\lambda\vdash n$ 的标准 Young 盘 $T$，以 $\{T\}$ 为基。相邻对换 $s_i=(i,i+1)$ 的作用为：
\begin{enumerate}[label=(\roman*),leftmargin=3em]
\item 若 $s_iT$ 仍为标准盘，且 $d=d_T(i,i+1)\neq0$：
\begin{equation}
\rho_\lambda(s_i)\ket{T}
=\frac{1}{d}\ket{T}+\sqrt{1-\frac{1}{d^2}}\,\ket{s_iT}.
\label{eq:young-yamanouchi-matrix}
\end{equation}
\item 若 $s_iT$ 不是标准盘（即交换 $i,i+1$ 后违反行或列递增）：
\begin{equation}
\rho_\lambda(s_i)\ket{T}=\frac{1}{d}\ket{T},
\label{eq:young-yamanouchi-diagonal}
\end{equation}
其中 $d=d_T(i,i+1)$。
\end{enumerate}
这组基是正交的，且 $\rho_\lambda$ 是\emph{实正交表示}。
\label{thm:young-yamanouchi}
\end{theorem}

\begin{proof}
在群代数 $\CC[S_n]$ 中引入 Jucys--Murphy 算子
\[
X_k=\sum_{r<k}(r\,k),\qquad k=1,\dots,n.
\]
它们彼此对易，并在标准盘基上对角化：$X_k\ket T=q_T(k)\ket T$。群代数恒等式
\[
s_iX_{i+1}-X_is_i=I,
\qquad
s_iX_k=X_ks_i\quad(k\ne i,i+1)
\]
说明 $s_i\ket T$ 只能落在 $\ket T$ 与 $\ket{s_iT}$ 张成的空间中。对第一式取 $\bra T$ 矩阵元，得到
\[
\bigl(q_T(i+1)-q_T(i)\bigr)\braket{T}{s_i}{T}=1,
\]
故对角元为 $1/d$。若 $s_iT$ 标准，则 $s_i^2=I$ 与实正交性把二维块唯一确定为
\[
\begin{pmatrix}
1/d&\sqrt{1-d^{-2}}\\
\sqrt{1-d^{-2}}&-1/d
\end{pmatrix}.
\]
若 $s_iT$ 不标准，则 $i,i+1$ 同行或同列，分别有 $d=1$ 或 $d=-1$，于是作用分别为 $+1$ 或 $-1$。由于这些算子就是 Specht 模中相邻对换的作用，Coxeter 关系随群代数表示自动成立。
\end{proof}


\par\medskip
\par\medskip

\begin{example}[\texorpdfstring{$S_3$ 的 {[2,1]} Young--Yamanouchi 表示}{S3 的 [2,1] Young-Yamanouchi 表示}]
形状 $[2,1]$ 有两个标准盘：
\[
T_1=\begin{array}{|c|c|}\hline1&2\\\hline3\\\cline{1-1}\end{array},
\qquad
T_2=\begin{array}{|c|c|}\hline1&3\\\hline2\\\cline{1-1}\end{array}.
\]
在 $T_1$ 中，$1,2$ 同行；在 $T_2$ 中，$1,2$ 同列。因此 $s_1$ 的本征值分别为 $+1,-1$。又有 $s_2T_1=T_2$，且 $d_{T_1}(2,3)=-2$、$d_{T_2}(2,3)=2$。所以在基 $(T_1,T_2)$ 中
\[
\rho(s_1)=\begin{pmatrix}1&0\\0&-1\end{pmatrix},
\qquad
\rho(s_2)=\begin{pmatrix}-1/2&\sqrt3/2\\\sqrt3/2&1/2\end{pmatrix}.
\]
直接相乘可验证 $\rho(s_i)^2=I$ 与 $\rho(s_1)\rho(s_2)\rho(s_1)=\rho(s_2)\rho(s_1)\rho(s_2)$。这就是 $S_3\cong D_3$ 的二维不可约表示。
\end{example}

\begin{example}[\texorpdfstring{$S_4$ 的 {[3,1]} Young--Yamanouchi 表示}{S4 的 [3,1] Young-Yamanouchi 表示}]
形状 $[3,1]$ 有三个标准盘：
\[
T_1=\begin{array}{|c|c|c|}\hline1&2&3\\\hline4\\\cline{1-3}\end{array},\quad
T_2=\begin{array}{|c|c|c|}\hline1&2&4\\\hline3\\\cline{1-3}\end{array},\quad
T_3=\begin{array}{|c|c|c|}\hline1&3&4\\\hline2\\\cline{1-3}\end{array}.
\]
对 $s_1=(12)$，$T_1,T_2$ 中 $1,2$ 同行，$T_3$ 中二者同列。对 $s_2=(23)$，$T_1$ 给出本征值 $+1$，而 $T_2,T_3$ 形成轴向距离为 $\mp2$ 的二维块。对 $s_3=(34)$，$T_1,T_2$ 形成轴向距离为 $\mp3$ 的二维块，$T_3$ 给出本征值 $+1$。因此
\begin{align*}
\rho(s_1)&=\operatorname{diag}(1,1,-1),\\
\rho(s_2)&=\begin{pmatrix}1&0&0\\0&-1/2&\sqrt3/2\\0&\sqrt3/2&1/2\end{pmatrix},\\
\rho(s_3)&=\begin{pmatrix}-1/3&2\sqrt2/3&0\\2\sqrt2/3&1/3&0\\0&0&1\end{pmatrix}.
\end{align*}
三个矩阵都是实正交的，满足 $S_4$ 的 Coxeter 关系，并生成三维不可约表示 $[3,1]$。
\end{example}

Young--Yamanouchi 表示的关键优势是：所有矩阵元素都是\emph{有理数或其平方根}，不需要复数运算；基向量由标准盘直接标记，不需要求解特征方程；相邻对换矩阵只在一个二维（或一维）子空间上非平凡，矩阵极其稀疏。这使得它在计算机实现中极为高效。

\section{低阶置换群的不可约表示}
\label{sec:low-sn-irreps}

抽象的 Young 方法只有在低阶例子中真正算一遍，才能看清它与\chapref{chap:representation-theory}的一般表示理论如何接上。这里重点算 $S_3$，再把 $S_4$ 的结构列出作为后续多电子问题的参考。

$S_3$ 的三个分拆是
\[
[3],\qquad[2,1],\qquad[1,1,1].
\]
钩长公式给维数
\[
1,\qquad2,\qquad1.
\]
$[3]$ 的行对称化意味着所有置换都作用为 $+1$，所以是平凡表示；$[1^3]$ 的列反对称化意味着置换按奇偶性作用，所以是符号表示。剩余 $[2,1]$ 给二维标准表示。

一个更透明的二维实现来自三维置换表示。$S_3$ 置换 $\CC^3$ 的标准基，向量
\[
\vb u_0=(1,1,1)
\]
张成平凡子表示，而平面
\[
W=\{(z_1,z_2,z_3)\mid z_1+z_2+z_3=0\}
\]
承载二维不可约表示。取正交基
\[
\vb e_1=\frac1{\sqrt2}(1,-1,0),
\qquad
\vb e_2=\frac1{\sqrt6}(1,1,-2),
\]
则可以直接算出
\[
D_{[2,1]}((12))=
\begin{pmatrix}-1&0\\0&1\end{pmatrix},
\]
以及三循环
\[
D_{[2,1]}((123))=
\begin{pmatrix}
-\tfrac12&-\tfrac{\sqrt3}{2}\\
\tfrac{\sqrt3}{2}&-\tfrac12
\end{pmatrix}.
\]
它恰好是平面内正三角形的二维几何表示。其特征标为
\[
\chi_{[2,1]}=(2,0,-1)
\]
对应恒等类、换位类与三循环类。

\begin{workedexample}
把置换
\[
\pi=\begin{pmatrix}
1&2&3&4&5&6&7\\
3&5&4&7&2&1&6
\end{pmatrix}
\]
分解为不交循环，并求其阶与奇偶性。

\textbf{解.} 从 $1$ 出发，$1\to3\to4\to7\to6\to1$，得到 $(13476)$；$2\to5\to2$，得到 $(25)$。故
\[
\pi=(13476)(25).
\]
阶为 $\operatorname{lcm}(5,2)=10$。五循环需要 $4$ 个对换，二循环需要 $1$ 个，总计奇数个，因此 $\operatorname{sgn}(\pi)=-1$。
\end{workedexample}


要对任意循环类型直接计算特征标，可以逐个移去对应循环长度的边钩，
避免构造完整表示矩阵。移去的图形必须连通，且符号由它跨过的行数决定。

\begin{definition}[边钩与高度]
若斜 Young 图 $\lambda/\nu$ 边连通，且不含 $2\times2$ 方块，
就称它为\textbf{边钩}。其长度是方格数，高度定义为占据的行数减一：
\[
h(\lambda/\nu)=\#\{\text{边钩占据的行}\}-1.
\]
这里 $\nu$ 是移去边钩后仍构成 Young 图的剩余形状。
\end{definition}

\begin{proposition}[Murnaghan--Nakayama 规则]
对 $\lambda\vdash n$ 与循环类型 $\mu=(\mu_1,\mu_2,\dots)\vdash n$，
\[
\chi^\lambda(\mu)=\sum_{\nu}(-1)^{h(\lambda/\nu)}\,
\chi^\nu(\mu_2,\mu_3,\dots),
\]
其中 $\nu$ 遍历满足 $\lambda/\nu$ 为长度 $\mu_1$ 的边钩的剩余形状。
依次移去所有循环长度，空图的特征标约定为 $1$；
若某一步无可移去的边钩，该项贡献为零。
\end{proposition}

\begin{example}[{手算 $\chi^{[2,1]}(3)$ 检验上面的 $(2,0,-1)$}]
循环类型 $(3)$ 意味着删一个长度 $3$ 的边钩。
$[2,1]$ 共三格，整图连通且没有 $2\times2$ 方块，故可以一次移去。
它占两行，高度为 $h=2-1=1$，符号为 $(-1)^1=-1$。
删完后图为空，递归终止贡献 $1$。于是
\[
\chi^{[2,1]}(3)=(-1)^{2-1}\cdot 1=-1,
\]
与几何表示在三循环 $(123)$ 处的特征标 $-1$ 一致。
平凡表示 $[3]$ 的边钩只占一行，高度 $0$，故特征标为 $+1$；
符号表示 $[1^3]$ 的边钩占三行，高度 $2$，故特征标也为 $+1$。
三者合起来正是 $S_3$ 特征标表的第三列 $(1,-1,1)$。
\end{example}

这里出现了贯穿全书的三重同构
\[
S_3\cong D_3\cong C_{3v},
\qquad
(123)\leftrightarrow r\leftrightarrow C_3,
\qquad
(12)\leftrightarrow s\leftrightarrow\sigma_v.
\]
\chapref{chap:group-basics} 从正三角形对称操作出发系统分析了 $D_3$ 的乘法表、共轭类、子群格和半直积结构；\chapref{chap:representation-theory} 算出了 $D_3$ 的二维不可约表示 $E$ 的显式矩阵\cref{eq:d3-E-matrices} 和特征标表\cref{tab:d3-character-table}；\chapref{chap:point-space-groups} 把同一个群实现为氨分子点群 $C_{3v}$。现在从置换群 $S_3$ 和 Young 图 $[2,1]$ 出发又到达了同一组矩阵。\emph{四个视角——几何对称、抽象群、分子点群、置换表示——给出同一个不可约表示}，这正是群论统一性的体现：同一个数学对象在不同物理语境中反复出现，每次都带着不同的物理含义但共享相同的代数骨架。

把同样的方法推进到 $S_4$，可以看到 Young 图怎样与熟悉的几何表示对应。

$S_4$ 的五个不可约表示维数为
\[
1,\ 3,\ 2,\ 3,\ 1.
\]
其中 $[4]$ 是平凡表示，$[1^4]$ 是符号表示；$[3,1]$ 是四维置换表示去掉平凡方向后的三维标准表示；$[2,1,1]$ 等于 $[3,1]$ 与符号表示张量积；$[2,2]$ 是二维不可约表示。

这与正四面体的置换图像有直接联系：$S_4$ 作用于四个顶点，而纯转动四面体群 $T\cong A_4$ 只包含偶置换。\chapref{chap:point-space-groups}中“正四面体的纯转动群为何只有 $12$ 个元素”现在可以用一个非常简洁的群论句子重述：
\[
T\cong A_4=\operatorname{Ker}(\operatorname{sgn}:S_4\to\{\pm1\}),
\qquad |A_4|=12.
\]
因此\chapref{chap:point-space-groups}的几何点群与本章的置换群并不是两个互不相关的专题，它们通过“群对几何顶点集合的作用”自然连接起来。

